
%\vspace{-1mm}

Bu çalışmada, güvenli e-posta iletişimini sağlamak amacıyla gizlilik, bütünlük, kimlik doğrulama ve inkar edememe özelliklerini bir arada sunan hibrit bir kriptografik sistem tasarlanmış ve gerçekleştirilmiştir. Sistem; Alice (gönderici) ile Bob (alıcı) iletişim senaryosu temelinde ele alınmış, her iki tarafın kriptografik işlemleri adım adım görselleştiren bir masaüstü uygulama arayüzüyle desteklenmiştir.

Önerilen sistemde mesaj bütünlüğü ve kimlik doğrulaması için SHA-256 özet fonksiyonu ile RSA-2048 PSS (Probabilistic Signature Scheme) dijital imza algoritması kullanılmaktadır. Mesaj gizliliği ise simetrik şifreleme algoritması AES-256-GCM (Galois/Counter Mode) ile sağlanmakta; simetrik oturum anahtarı RSA-2048 OAEP dolgusu kullanılarak asimetrik olarak şifrelenmektedir. Bu hibrit yaklaşım sayesinde asimetrik kriptografinin güvenlik avantajları ile simetrik kriptografinin performans avantajları bir arada elde edilmektedir.

Gerçekleştirilen PyQt6 tabanlı masaüstü uygulama; AES-256, RSA-2048 ve SHA-256 algoritmalarının iç işleyişini 14 turlu AES durum matrisi, Genişletilmiş Öklid Algoritması adımları ve SHA-256 sıkıştırma devresi animasyonlarıyla pedagojik biçimde sunmaktadır. Uygulama 26 birim testi ile doğrulanmış olup NIST FIPS 180-4, NIST SP 800-38D ve RFC 8017 standartlarına uygun kriptografik primitifler kullanılmaktadır.

\vspace{3mm}

 \begin{tabular}{ll}
{\textbf{Anahtar Kelimeler:}} & 
{Hibrit Şifreleme, AES-256-GCM, RSA-2048, SHA-256, Dijital İmza, E-posta Güvenliği.} 	 
 \end{tabular} 

